\subsection{E-graphs and Equality Saturation}
An e-graph (equivalence graph) is a data structure that stores equivalence relations in a compact form to extract the most optimal expression later.
The e-graph consists of nodes, called e-nodes in the e-graph terminology. 
Equivalent e-nodes are grouped together in e-classes where all nodes in a single e-class are equivalent. 
E-nodes contain expressions which can be functions or variables.
The expression is constructed by drawing an arrow from an e-node representing the function to the e-classes containing its arguments. 
For instance, in Figure \ref{fig_egg_expr},  the node $/$ represents the expression $(*$ e-node$)/2$. Then the $*$ e-node represents the expression $a*2$. 
Therefore, The $/$ node represents the expression $(a*2)/2$. 
It is important to note that arrows are drawn from e-nodes to e-classes, not to other e-nodes.
The user also defines the criteria for choosing the optimal expression in the form of a cost function that is calculated while traversing the e-graph.

The construction and usage of the e-graph consists of two stages: building the e-graph based on rewrite rules until equality saturation, and extraction phase.
Rewrite rules are defined by the user to specify a list of equivalent expressions.
The first stage is to build the e-graph using the rules until no rewrite rule adds further information to the e-graph, in which case the e-graph has reached the equality saturation phase. 
For example, the expression $a*2$ is equivalent to $a <<1$.
In Figure \ref{fig_egg_div_eq}, the two expressions are represented using two separate graphs. 
When using e-graphs, the same information can be represented in a much more compact form as shown in Figure \ref{fig_egg_div_egraphs}. 
The second stage is meant to traverse the e-graph and extract the most optimal expression based on a cost function defined by the user. 
The good thing about e-graphs is that it allows for non-destructive rewriting, which means it does not forget old rules, but rather only adding and growing the e-graph until saturation. 
Moreover, given the compactness nature of e-graphs, it can represent thousands of nodes and classes without a problem.  

\begin{figure}[H]
    \centering
    \begin{subfigure}{0.15\textwidth}
        \centering
        \includegraphics[width=\textwidth]{sections/figures/egg/egg_div_expr.pdf}
        \caption{Divide Expression using E-graph}
        \label{fig_egg_expr}
    \end{subfigure}
    \hfill
    \begin{subfigure}{0.3\textwidth}
        \centering
        \includegraphics[width=\textwidth]{sections/figures/egg/egg_div_mul_eq.pdf}
        \caption{Equivalent Multiplication and Shift Left Operations as two separate graphs}
        \label{fig_egg_div_eq}
    \end{subfigure}
    \hfill

    \vspace{1em} % Vertical space between rows
    
    \begin{subfigure}{0.15\textwidth}
        \centering
        \includegraphics[width=\textwidth]{sections/figures/egg/egg_div_mul_egraph.pdf}
        \caption{Equivalent Multiplication and Shift Left Operations in a Single E-graph}
        \label{fig_egg_div_egraphs}
    \end{subfigure}
       \caption{Rewrite Rules in E-graphs}
       \label{fig_egg_rewrite}
\end{figure}

Authors in \cite{panchekhaEggFastExtensible2021} introduced egg, a Rust library which implements e-graphs and equality saturation principles for easier integrations which made it a good fit for many different applications. 
For instance, egg can be used to optimize floating point's numerical stability \cite{tatlockEggFloatingPoint2015}, and optimize linear algebra implementation \cite{danEggLinearAlgebra2020}. 
Moreover, authors in \cite{draneAutomaticDatapathOptimization2022} designed rewrite rules and used egg to automate datapath RTL optimization.
User only defines the functions, which is called a language in egg term, and defines the rewrite rules that needs to be performed on certain expressions. 
For more flexibility, egg was integrated with Datalog principles so developers can create egg programs without the need to the Rust library \cite{zhangEGGLog2023}. 
For instance, authors in \cite{shanabrookEgglogPythonPythonic2024} created python-egglog making egg even easier to write by implementing a Python abstraction layer that interacts with the Rust library. 
Additionally, python-egglog gives users more flexibility to write a pure egglog program and interacts with it through python interface and added more functionality to create egglog programs in pythonic syntactic and functional style. 